%%%%%%%%%%%%%%%%%%%%%%%%%%%%%%%%%%%%%%%%%%%%%%%%%%%%%%%%%%%%
%% Chapter 8 - Conclusion and Future Work
%% Summarise what was established, answer the research questions,
%% identify next steps. No new experimental evidence.
%%%%%%%%%%%%%%%%%%%%%%%%%%%%%%%%%%%%%%%%%%%%%%%%%%%%%%%%%%%%
\chapter{Conclusion and Future Work}
\label{chap:conclusion}
\fancyhead[RE]{\headfont\nouppercase{\leftmark}}

This thesis asked how hydraulic actuator dynamics learned from high-fidelity
simulation can be reduced, transferred into a computationally efficient robotic
simulation and used for reinforcement-learning control, and under what
conditions online adaptation adds a further increment to the robustness those
learned models already provide. The six research questions of
Section~\ref{sec:intro:contributions} are answered in the order they were
asked, which is also the order in which the answers depend on one another.

\section{Summary of Findings}

The findings are stated in the order of the research questions. The first four
concern how hydraulic actuation is learned, transferred and controlled, and the
last two concern robustness and online adaptation. Figure~\ref{fig:concl:evidence-chain}
makes the dependency between those experiments explicit.

\begin{figure}[htbp]
  \centering
  \begin{tikzpicture}[
    evidence/.style={draw=black!65, rounded corners=3pt, line width=.7pt,
      minimum height=14mm, text width=.82\textwidth, align=left,
      inner xsep=7pt, inner ysep=5pt, font=\small, fill=white},
    flow/.style={-{Latex[length=2.4mm]}, line width=.9pt, draw=black!65}
  ]
    \node[evidence, draw=cbblue, line width=1.1pt] (models) {\textbf{1. Representation and transfer}\quad
      Simscape identification, open-loop excitation, command ablation and
      recursive rollout establish usable ID--FD and ROM actuator paths. The
      ID--FD path transfers frozen; the ROM receives bounded deployment
      calibration.};
    \node[evidence, draw=cbgreen, line width=1.1pt, below=6mm of models] (control) {\textbf{2. Control and robustness}\quad
      Teacher distillation, the matched R--IF comparison, 33 evaluation cells
      and plant-factor sweeps establish policy-only control and the tested
      robustness and authority margin.};
    \node[evidence, draw=cborange, line width=1.1pt, below=6mm of control] (adapt) {\textbf{3. Adaptation mechanisms}\quad
      The offline AAM tests observability and is usable at the model level; the
      latent-delivery experiment exposes saturation at that interface; direct-valve
      AF trials then separate frozen-base, residual-capacity, history and
      plant-dependent effects.};
    \draw[flow] (models) -- (control);
    \draw[flow] (control) -- (adapt);
  \end{tikzpicture}
  \caption{Experimental evidence chain used to answer the research questions.
  Each stage constrains the claim made at the next stage; the final adaptation
  claim is interpreted only after transfer and robustness have been established.}
  \label{fig:concl:evidence-chain}
\end{figure}

\textbf{The inverse and forward representation identified in Simscape transfers
to the deployment simulation without retraining (RQ1, RQ2).} The inverse model
was trained on Simscape data alone and evaluated on all $118$ IsaacLab
deployment runs, drawn entirely outside that training distribution, where its
teacher-forced one-step extension error is millimetre-scale and nearly flat
across domains: $\unit[2.36]{mm}$ on the combined runs against $\unit[3.36]{mm}$
on nominal. The claim holds only in that one-step regime; closing the loop on
its own predictions diverges in all $118$ runs
(Section~\ref{sec:models:id-trust}). The reduced-order model received a small
bounded calibration against the deployment plant.

\textbf{Predictive accuracy must be accompanied by a control-authority test
(RQ1).} A learned actuator must be tested for dependence on the command, not
only for prediction error. With the accepted excitation design, removing the
valve input increases rollout error by a factor of $1.81$; before it, by $0.94$.
This is a general validation requirement for learned models used in control.

\textbf{Two actuator representations both support closed-loop control, and they
do not produce the same controller (RQ3, RQ5).} The command-conditioned
reduced-order model reaches $R^2 = 0.998$ on the internal valve state and $0.97$
on extension, and runs inside GPU-parallel training. The inverse--forward chain
composes two separately learned models in series, which could have compounded
their errors and did not. The two induce different control strategies from the
same policy interface. On the
recorded traces of one held-out mesh, \ppoROM{} spends the least command and
reverses about $40$ times per second against $10.6$ for \ppoIDFD{}, which
converts greater actuator effort into lower lateral error
(Table~\ref{tab:eval:effort}). Over the $33$ paired cells the median cell
favours \ppoIDFD{} on lateral error by $\unit[9.3]{mm}$ with the interval
excluding zero and favours \ppoROM{} on cut height by $\unit[1.7]{mm}$, while
neither mean difference is resolved and the mean changes sign if the double
S-curve family is removed (Table~\ref{tab:eval:paired-rif}). The choice of
representation therefore affects the resulting controller and is not an
implementation detail, but the size of that effect is not established.

\ppoROM{} and \ppoIDFD{} match on recorded training fields and differ only in
the actuator path. AF uses a different seed and training configuration and is
therefore evaluated through the paired within-stack comparison of
Section~\ref{sec:eval:attribution}, not ranked against them. With one R and one
IF training run, their seed variance remains unknown.

\textbf{A learned policy can replace the classical controller at runtime
(RQ4).} The damped-least-squares controller supplies privileged demonstrations
during training and is then removed; the distilled policy emits desired cylinder
motion directly into the learned actuator and is the entire deployed controller.
Over $33$ evaluation cells \ppoIDFD{} holds a mean lateral error of
$\unit[70.3]{mm}$, a median of $\unit[28.1]{mm}$ and $95.29\,\%$ ground contact
(Table~\ref{tab:eval:tracking}). The deployed policies remain anchored to the
demonstrated behaviour because the reward is an imperfect proxy for the
evaluation gates.

\textbf{The robust stacks retain a broad margin against mechanical variation
and authority loss without online adaptation (RQ2, RQ5).} On held-out meshes,
eleven trajectory families and sixteen reserved factor configurations,
\ppoIDFD{} meets the lateral criterion in $32$ of $33$ cells at a median of
$\unit[0.0281]{m}$. Its frozen inverse model transfers to all $118$ deployment
runs with millimetre-scale one-step error.

The factor sweep supports the tested mechanical variation, but not unbounded
hydraulic variation. On the inverse and forward path, hydraulic factors are not
applied; on \ppoROM{} they remain inside the identification range
(Section~\ref{sec:eval:delivered-stress}). Reducing realised actuator motion to
one per cent of nominal did not trigger the safety or ground-contact gates.
The result therefore establishes an authority margin, not a hard robustness
boundary (Section~\ref{sec:eval:boundary-design}).

\textbf{Online plant context is usable, and whether it reaches the closed loop
is set by the correction interface (RQ6).} The explicit mechanism estimates a
plant latent from recent interaction history. Its latent-conditioned adapters
beat a zero-context control on the adapter objective on all five seeds, and the
same correction changes no closed-loop metric, because it is added to an inverse
command already resting against its limit for most of the scored window. Moving
the correction onto the valve path restores the authority the residual needed.
At the pre-registered endpoint the history-conditioned residual succeeds on four
of five seeds; the fixed final checkpoint succeeds on three. An equal-capacity
zero-history residual also beats the frozen base on four seeds, so residual
capacity explains the median gain. Live history separates from that control on
one seed under each checkpoint rule. Restricting the encoder to plant-relevant
channels, the leading explanation for that absence, does not recover it: on a
plant where the frozen base degrades by two orders of magnitude across the mass
range, the primary comparison succeeds on two seeds of five and the median seed
fails (Section~\ref{sec:eval:starved}). Capacity alone beats the frozen base on
four of five there, one more than the history-conditioned arm. The cutter-mass study resolves a
load-dependent lateral component but compares against the frozen base, so it
cannot distinguish capacity from history. The supported method-level claim is
therefore the bounded residual; the history-specific increment is seed-dependent.

Table~\ref{tab:concl:adaptation-evidence} separates the experiments behind that
conclusion, so the gain of an enlarged residual is not read as evidence that
live history caused it.

\begin{table}[htbp]
  \centering
  \caption{Experimental decomposition of the adaptation claim. The rows answer
  different questions at different levels of the control chain and are not
  interchangeable evidence.}
  \label{tab:concl:adaptation-evidence}
  \small
  \setlength{\tabcolsep}{4pt}
  \begin{tabular}{@{}p{.20\textwidth}p{.27\textwidth}p{.43\textwidth}@{}}
    \toprule
    Experiment & Controlled comparison & Supported conclusion \\
    \midrule
    Latent adapter objective
      & estimated, privileged and zero latent on the adapter's own objective
      & Online plant context is usable at the model level: both latent arms beat the zero-context control on five of five seeds. \\
    Latent-delivery test
      & estimated, privileged and zero latent through the inverse-model correction path
      & Better plant context does not recover extension or velocity; saturation leaves insufficient correction authority at that interface. \\
    Primary seed table
      & history residual against its frozen base, paired within seed
      & The residual improves tracking on four of five seeds and the median seed passes; three of five at the final checkpoint. \\
    Capacity control
      & history-conditioned residual against an equal-capacity zero-history residual
      & Capacity alone reproduces most of that gain; live history separates from zeroed history on one seed of five at each checkpoint rule. \\
    Encoder restriction
      & history against zeroed history at equal capacity, encoder starved to plant-relevant channels
      & The missing history effect is not an artefact of a wide encoder input: two seeds of five succeed and the median seed fails. \\
    Cutter-mass dose and sham
      & plant perturbation against the corresponding sham condition
      & A plant-dependent lateral component is resolved, but starting mass and the later exchange remain coupled. \\
    \bottomrule
  \end{tabular}
\end{table}

\section{Design Implications}

The design implications follow the same sequence: actuator identification,
policy training and plant adaptation.

\subsection*{Identification and Transfer}

\textbf{Design the excitation for identifiability, and test the model by input
ablation.} Task-control data can make state and command nearly redundant.
Dedicated excitation makes command effects identifiable. A learned actuator
should therefore be accepted on both prediction accuracy and loss of performance
when the command input is removed.

\textbf{Embed known physics and bound learned corrections.} The reduced-order
model propagates pressure, force, velocity, and extension through explicit
updates. Its learned terms are bounded, preventing unphysical state changes
inside the training loop. This structure also represents the pressure dependence
of valve authority directly.

\textbf{Prefer a frozen model with a bounded, gated correction to retraining.}
The inverse model transferred at millimetre scale and remained frozen. A
$\unit[3]{mm}$ correction absorbed the deployment offset, vanished at zero
command and preserved the original data through replay, confining
deployment-specific behaviour to a bounded component.

\subsection*{Control Methodology}

\textbf{Constrain policy updates when the reward is not the evaluation
criterion.} The physical evaluation gates and the reward are not identical. A
forward-KL anchor to a verified reference policy bounds policy drift, and the
anchor advances only through checkpoints that satisfy the gates. Reward terms
should represent quantities that are also treated as evaluation constraints.

\textbf{Where no baseline is inherited, the action contract is the safety
argument.} Zero action leaves the arm stationary. A rate limit, filter, travel
shield, command scale and per-axis anti-windup cap bound its authority. The cap
must be resolved per axis because the value that recovers boom offsets damages
head alignment on cylinder~4 (Section~\ref{sec:ctrl:ppo}).

\subsection*{Plant Adaptation}

\textbf{Measure identifiability and correction authority before training an
estimator.} The identifiability envelope determines whether plant identity is
present in the data. The delivered-stress check determines whether a
perturbation reaches the plant and whether a correction can influence the
command. Both tests should precede estimator training.

Filtering may be needed when a changing estimate enters a memoryless policy,
but its time constant remains to be determined on the deployed stack.

\section{Contributions}

Section~\ref{sec:intro:contributions} stated six contributions. Their standing
at the end of the work is as follows, with the supporting evidence as summarised
above rather than repeated.

\textbf{Transferable hydraulic actuator learning} is delivered through two
Simscape-identified representations realised in IsaacLab; the inverse model
transfers without retraining. The \textbf{reduced-order actuator model} is a
bounded, differentiable surrogate suitable for GPU-parallel training and serves
as the calibrated deployment plant. Its clamped constants are not identified
physical parameters (Section~\ref{sec:models:rom}).

The \textbf{control-relevant identification and validation} procedure combines
input ablation, which exposed a forward model without command authority, with an
identifiability audit of the excitation envelope.

\textbf{Policy-only control through learned hydraulics} is demonstrated. After
the classical teacher is removed, \ppoIDFD{} holds a mean lateral error of
\unit[70]{mm} and a median of \unit[28]{mm} over $33$ held-out cells at
$95\,\%$ ground contact (Section~\ref{sec:eval:tracking}). Worst-cell
aggregation still misses two gates; that bound is stated in
Section~\ref{sec:concl:limitations}.

The \textbf{decision procedure for online adaptation} separates residual plant
sensitivity, observability, correction authority and residual capacity before
attributing a gain to adaptation. It also exposes when the optimiser trades away
a gated quantity that the reward does not price.

\textbf{Where online adaptation is effective} is answered at two levels. Plant
context is informative at the model level but cannot affect the saturated
inverse command. A bounded valve-path residual retains authority and improves
the median training run, although the capacity-matched control reproduces most
of the gain. Live history contributes only on seed-specific cases, and a second
campaign that starves the encoder to plant-relevant channels returns two of five
seeds, eliminating the encoder as the explanation. The cutter-mass study resolves
a plant-dependent lateral component without isolating the transient mass
exchange.

Two enablers are recorded: the causal estimator and bounded adaptation path use
no privileged or future information and revert safely on invalid input. Five of
the eight promotion gates separate the matched pair
(Section~\ref{sec:sys:gates}).

\section{Limitations}
\label{sec:concl:limitations}
\label{sec:conclusion:limitations}

\textbf{Scope of plant validation.} The evidence is simulation-to-simulation,
and the deployment plant is itself a learned model. No claim here is a claim
about hardware.

\textbf{Scope of hydraulic observability.} The differential pressure the estimator consumes is produced by the actuator
model rather than measured, so transferring the model-level result to hardware
requires instrumentation the FME~500 does not carry. An observation contract
remains unresolved alongside this: every result uses the pressure-inclusive row,
while the stricter deployable alternative excludes measured differential
pressure. Since the estimator's advantage is concentrated in that channel, the
deployable variant should be expected to perform worse, and by how much is not
established.

\textbf{Scope of controller comparison.} On worst-mesh, worst-family statistics
both stacks fail the lateral gate and \ppoIDFD{} does not satisfy the
head-alignment gate under worst-case aggregation
(Section~\ref{sec:eval:tracking}), so neither stack satisfies the whole
eight-gate evaluation contract. The closed-loop comparisons resolve differences to a few per cent over $33$
paired conditions and do not establish equivalence. The available plant
variation bounds the disturbance rather than the architecture, and more tightly
than the sweep's labels suggest. A stack has been shown to tolerate what
actually reached it, which on the inverse and forward plant is a mechanical
perturbation only, plus a generic reduction of realised actuator motion. Median
lateral error first reaches $1.5$ times its full-authority value at
$\kappa = 0.20$ for \ppoIDFD{} and $\kappa = 0.10$ for \ppoROM{}, and no stack
reaches the pre-declared breakdown gate down to $\kappa = 0.01$
(Section~\ref{sec:eval:boundary-design}). That is an authority margin, not a
demonstration of hydraulic robustness. The three
stacks are matched on evaluation but not on training budget, so the comparison
describes three deployed controllers rather than three equally tuned training
procedures.

\textbf{Scope of training variation.} Policy-training variation is not quantified for
the matched R and IF pair, and it bounds
two claims differently. R and IF are each represented by one training run, so
the $33$ paired cells measure terrain and trajectory variation and say nothing
about how far a second run of either would move; the actuator-representation
result rests on that assumption. The adaptation result includes five training
seeds and meets the primary criterion on four, so it establishes a median
improvement and not a uniform one. The defensible claim is an
architecture-level correction on the identified valve path, together with a
plant-dependent lateral term under the tested cutter-mass doses.

\textbf{Scope of the adaptation evidence.} The commanded step cannot finish faster than $417$\,ms at the command-rate limit,
while the plant answers in about $117$\,ms. That loop is a hard case for
adaptation understood as a reaction to a single event and favours continuous
conditioning on the operating point.
Appendix~\ref{app:pressure-signals} is consistent with the second reading:
the adaptive advantage is already present before the load change.
Section~\ref{sec:eval:dose} shows a cutter-mass-dependent lateral term relative
to the $1.0$ condition, which the coupled dose design does not resolve into a
transient adaptation response.

\section{Future Work}
\label{sec:concl:future}

The evidence in this thesis is simulation-to-simulation, and the three
experiments that would most change what it can claim are also in simulation.
They are listed first for that reason.

\textbf{Fit the encoder by regression rather than through the policy gradient.}
Restricting the encoder input has been done and did not recover the effect
(Section~\ref{sec:eval:starved}), which leaves the training signal rather than
the input width as the remaining candidate. Rapid motor adaptation never trains
its encoder through the policy gradient, fitting it onto a low-dimensional
privileged vector and freezing it before the residual head is
trained~\cite{kumar2021rma}. That is the one variant this campaign's
pre-registration admits, and it separates a representation problem from the
absence of a usable effect.

\textbf{Train more than one seed per representation.} R and IF are one training
run each, so the cross-stack comparison of Section~\ref{sec:eval:tracking}
describes two deployed controllers and cannot separate the representation from
the seed. Three or more runs per stack, scored on the same $33$ cells, would
convert that observation into a result at moderate cost.

\textbf{Ablate reference-anchored PPO.} Section~\ref{sec:ctrl:finding}
establishes that unconstrained fine-tuning fails; it does not establish that
this particular construction is what fixed it. A three-arm comparison against
plain PPO with early stopping and against the un-fine-tuned distilled policy
would separate the block ratchet from the constraint it sits inside.

\textbf{Then the machine.} Hardware evaluation begins by adding chamber-pressure
and extension sensing and a rod load pin to the FME~500. The open-loop excitation
schedule and identifiability audit of Chapter~\ref{chap:data} can then be applied
to machine data. Any correction should remain bounded, retain the identified
hydraulic parameters and original data for replay, and pass input ablation. The
reference-anchored policy provides the verified starting checkpoint.

The central design lesson is conditional. Robust training absorbs a large part
of the plant variation before online adaptation is introduced. A remaining
correction is useful only when the variation is observable, the fixed controller
leaves a deficit to recover, and the command path retains authority. Apparent
adaptation gains must also be separated from gains due to residual capacity.
Evaluating these conditions together provides a defensible basis for deciding
when online adaptation is justified on a hydraulic manipulator.
